% Ad Familiares 14.5 --- headnote
% ad-familiares-14-05, 16 October 50 BC, Athens

\textit{Cicero to Terentia, written from Athens on the seventeenth
day before the Kalends of November (16 October) 50 BC (the manuscript
dateline: \textit{Scr. Athenis a. d. xvii K. Novemb. a. 704 (50)}).
The proconsul is on his way home from Cilicia, ship just landed, the
freedman Acastus on the quay with letters from Italy. The opening
formula --- ``if you and Tullia are well, then I and our dearest
Cicero are well'' --- is the standard salutation of every letter to
Terentia, and is here so framed because Tullia's health is what
Cicero is most anxious about; he asks her, ``so far as the state of
your health allows,'' to come to meet him.}

\textit{The rest is hurried business: the bad crossing, the Precius
inheritance to be managed by Atticus (Pomponius) or by Camillus if
he is not available, the household plans for the auction. Underneath
is the public news the friends' letters have just delivered ---
``the matter is tending toward arms'' --- which means his return will
be a political crisis, not the relief of having put Cilicia behind
him. He hopes to be in Italy by the Ides of November. The civil war
is two months off.}
